% =====================================================================
\chapter{Regresión lineal simple y correlación}
\label{tema:6}
% =====================================================================

La regresión estudia cómo cambia en promedio una variable respuesta $y$ cuando cambia una variable explicativa $x$. Este tema deduce la recta de mínimos cuadrados, estudia los residuos y la descomposición de la variabilidad que da lugar al coeficiente de determinación $R^{2}$, muestra cómo las transformaciones logarítmicas permiten ajustar leyes de potencia y termina con la distinción entre correlación y causalidad. Se sigue el capítulo~12 de OpenStax~\cite{openstax} (§12.3 «The Regression Equation»), el capítulo~14 de Rice~\cite{rice}, «Linear Least Squares», el capítulo~13 de Wasserman~\cite{wasserman}, «Linear and Logistic Regression», y la lectura~26 del curso 18.05 del MIT~\cite{mit1805}; para la causalidad, el capítulo~16 de Wasserman, «Causal Inference».

\section{La recta de mínimos cuadrados}

Se dispone de $n\ge2$ pares $(x_1,y_1),\dots,(x_n,y_n)$ con los $x_i$ no todos iguales. Se usan las sumas centradas
\[
S_{xx}=\sum_{i=1}^{n}(x_i-\xbar)^{2},\qquad S_{yy}=\sum_{i=1}^{n}(y_i-\ybar)^{2},\qquad S_{xy}=\sum_{i=1}^{n}(x_i-\xbar)(y_i-\ybar),
\]
de modo que $s_x^{2}=S_{xx}/(n-1)$, $s_{xy}=S_{xy}/(n-1)$ y $r=S_{xy}/\sqrt{S_{xx}S_{yy}}$ (\cref{def:t1-cov}). Por la \cref{prop:t1-abreviada}, $S_{xx}=\sum x_i^{2}-n\xbar^{2}$ y, análogamente, $S_{xy}=\sum x_iy_i-n\xbar\ybar$.

\begin{definicion}[Recta de mínimos cuadrados]\label{def:t6-mc}
Para cada recta $y=b_0+b_1x$, el \emph{valor ajustado} del dato $i$ es $\widehat y_i=b_0+b_1x_i$ y el \emph{residuo} es $e_i=y_i-\widehat y_i$. La \emph{recta de mínimos cuadrados} (o de regresión de $y$ sobre $x$) es la que minimiza la suma de cuadrados de los residuos
\[
\operatorname{SC}(b_0,b_1)=\sum_{i=1}^{n}(y_i-b_0-b_1x_i)^{2}.
\]
\end{definicion}

\begin{teorema}[Coeficientes de mínimos cuadrados]\label{teo:t6-mc}
Si $S_{xx}>0$, la función $\operatorname{SC}$ tiene un único mínimo, en
\[
\widehat b_1=\frac{S_{xy}}{S_{xx}}=r\,\frac{s_y}{s_x},\qquad\widehat b_0=\ybar-\widehat b_1\xbar,
\]
y el valor mínimo es $S_{yy}-S_{xy}^{2}/S_{xx}=S_{yy}(1-r^{2})$.
\end{teorema}

La condición $\widehat b_0=\ybar-\widehat b_1\xbar$ dice que la recta pasa por el punto de medias $(\xbar,\ybar)$. La pendiente tiene el signo de $r$: es positiva cuando la asociación es creciente.

\addtocounter{definicion}{1}%

\addtocounter{figure}{1}%

\section{Residuos y coeficiente de determinación}

\begin{proposicion}[Propiedades de los residuos]\label{prop:t6-residuos}
Para la recta de mínimos cuadrados, $\sum_ie_i=0$, $\sum_ix_ie_i=0$ y $\sum_i\widehat y_ie_i=0$.
\end{proposicion}

La demostración es el \cref{eje:t6-residuos}. Las dos primeras igualdades son las \emph{ecuaciones normales} (\cref{eje:t6-normales}): los residuos son ortogonales, como vectores de $\R^{n}$, al vector de unos y al vector de las $x_i$. Este es el punto de vista geométrico que el módulo~8 generaliza a varias variables explicativas.

\begin{teorema}[Descomposición de la variabilidad]\label{teo:t6-anova}
Sean $\operatorname{SC}_T=\sum_i(y_i-\ybar)^{2}=S_{yy}$ (suma de cuadrados total), $\operatorname{SC}_R=\sum_i(\widehat y_i-\ybar)^{2}$ (explicada por la regresión) y $\operatorname{SC}_E=\sum_ie_i^{2}$ (residual o de error). Entonces
\[
\operatorname{SC}_T=\operatorname{SC}_R+\operatorname{SC}_E.
\]
\end{teorema}

\begin{definicion}[Coeficiente de determinación]\label{def:t6-r2}
El \emph{coeficiente de determinación} es $R^{2}=\operatorname{SC}_R/\operatorname{SC}_T=1-\operatorname{SC}_E/\operatorname{SC}_T$, la fracción de la variabilidad de $y$ explicada por la recta. Por el \cref{teo:t6-anova}, $0\le R^{2}\le1$.
\end{definicion}

\begin{corolario}\label{cor:t6-r2}
En la regresión lineal simple, $R^{2}=r^{2}$.
\end{corolario}

En el \cref{ej:t6-recta}, $R^{2}=0{,}6$: la recta explica el $60\,\%$ de la variabilidad de $y$. La dispersión de los puntos alrededor de la recta se resume en el \emph{error típico residual} $s_e=\sqrt{\operatorname{SC}_E/(n-2)}$; se divide por $n-2$ porque se han estimado dos parámetros (compárese con la \cref{obs:t1-n-1}).

\section{Diagnóstico con residuos}

Un $R^{2}$ alto no garantiza que el modelo lineal sea adecuado. La herramienta básica de diagnóstico es el \emph{gráfico de residuos}, que representa $e_i$ frente a $x_i$ (o frente a $\widehat y_i$). Si el modelo es adecuado, los residuos se distribuyen sin estructura alrededor de cero. Algunas señales de alarma~\cite[§12.3]{openstax}:
\begin{enumerate}
  \item una curvatura sistemática (residuos positivos en los extremos y negativos en el centro, o al revés) indica que la relación no es lineal;
  \item una dispersión que crece con $x$ (forma de embudo) indica varianza no constante;
  \item un punto con un residuo muy grande es un posible atípico, y un punto con $x_i$ muy alejado de $\xbar$ puede ser muy \emph{influyente}: puede cambiar mucho la recta por sí solo.
\end{enumerate}
Además, la recta solo describe la relación en el rango de los $x_i$ observados: \emph{extrapolar} fuera de ese rango (\cref{eje:t6-buceo}) es arriesgado.

\begin{observacion}[El modelo estadístico]\label{obs:t6-modelo}
Si se supone que $y_i=\beta_0+\beta_1x_i+\varepsilon_i$ con errores $\varepsilon_i\sim\Nor(0,\sigma^{2})$ \iid, la log-verosimilitud de $(\beta_0,\beta_1)$ es una constante menos $\frac{1}{2\sigma^{2}}\operatorname{SC}(\beta_0,\beta_1)$; por tanto, los estimadores de mínimos cuadrados son los estimadores de máxima verosimilitud~\cite[cap.~13]{wasserman}, \cite[\emph{problem set}~6]{mit1805}. Con este modelo pueden construirse intervalos de confianza y contrastes para la pendiente (ampliación), que se estudian en~\cite[cap.~14]{rice}.
\end{observacion}

\section{Transformaciones: ajuste de leyes de potencia}

Muchas relaciones no lineales se vuelven lineales tras una transformación. Una \emph{ley de potencia} $y=a\,x^{b}$ con $a,x,y>0$ cumple
\[
\log y=\log a+b\log x,
\]
de modo que en escala log-log es una recta de pendiente $b$ y ordenada $\log a$. Para ajustarla se aplica mínimos cuadrados a los pares $(\log x_i,\log y_i)$ y se deshace la transformación: $\widehat a=10^{\widehat b_0}$ (o $e^{\widehat b_0}$ si se usan logaritmos naturales) y $\widehat b=\widehat b_1$. Análogamente, una relación exponencial $y=a\,e^{bx}$ se linealiza con $\ln y=\ln a+bx$. Obsérvese que así se minimizan los errores \emph{relativos} (en escala logarítmica), no los absolutos.

\addtocounter{definicion}{1}%

\section{Correlación y causalidad}

Una correlación entre $x$ e $y$ puede deberse a que $x$ influye en $y$, a que $y$ influye en $x$ (causalidad inversa), a que una tercera variable $z$ influye en ambas (\emph{variable de confusión}) o simplemente al azar. Los datos observacionales, por sí solos, no permiten distinguir estos casos~\cite[cap.~16]{wasserman}.

\addtocounter{definicion}{1}%

Un caso extremo de confusión es la \emph{paradoja de Simpson}: una asociación que se da en cada uno de los grupos de una población puede invertirse al agregar los grupos, cuando los tamaños de los grupos difieren entre las condiciones comparadas (\cref{eje:t6-simpson}).

La forma más fiable de establecer una relación causal es el \emph{experimento aleatorizado}: si el tratamiento (por ejemplo, el \emph{prompt} A o el B) se asigna al azar a las unidades (las preguntas), las variables de confusión se reparten en promedio por igual entre los grupos, y una diferencia sistemática en la respuesta solo puede atribuirse al tratamiento o al azar, cuya magnitud se cuantifica con los contrastes del \cref{tema:5}.

\addtocounter{definicion}{1}%

\addtocounter{definicion}{1}%

